\documentclass{article}
\usepackage[T1]{fontenc}
\usepackage{lmodern}
\usepackage{graphicx}
\usepackage{amsmath,amssymb}
\usepackage[a4paper,margin=2cm]{geometry}
\usepackage[absolute,overlay]{textpos}
\setlength{\TPHorizModule}{1mm}
\setlength{\TPVertModule}{1mm}
\usepackage[indonesian]{babel}
\usepackage{hyperref}
\setlength{\emergencystretch}{2em}
% unit: Halaman 58

\begin{document}

% === Halaman 58 (python/native) ===

58

• Kelompokkan “pesan” setiap berjarak 5, dimulai dari huruf cipherteks pertama, 

kedua, dan seterusnya. Setiap huruf pada posisi berjarak 5 pasti dienkripsi 

dengan huruf kunci yang sama.

   1234567891011…

 LJVBQSTNEZLQMEDLJVMAMPKAUFAVATLJVDAYYVNFJQLNPLJVHKVTRNFLJVCM 

LKETALJVHUYJVSFKRFTTWEFUXVHZNP


    Kelompok 

Pesan  



    Huruf paling sering muncul

  1  

LSLLMFLYJLVLLLYKWV 


L

   2  

JTQJPAJYQJTJKJJREH 


J

   3  

VNMVKVVVLVRVEVVFFZ 


V

   4  

BEEMAADNNHNCTHSTUN 


N

   5  

QZDAUTAFPKFMAUFTXP 


A

\end{document}
